%=========== ΛΥΣΗ - 22ο ΘΕΜΑ Α ===========
\providecommand{\theoriaref}[3]{%
\begin{center}
\fcolorbox{maincolor!70!black}{maincolor!8}{%
\begin{minipage}{.88\linewidth}
\textcolor{maincolor!80!black}{\kerkissans{\textbf{#1~\ref{#2}}}}\\[1mm]
#3
\end{minipage}}
\end{center}}
\providecommand{\orismosref}[1]{%
\theoriaref{Ορισμός}{#1}{Η απάντηση δίνεται από τον αντίστοιχο ορισμό.}}
\providecommand{\apodeixiref}[1]{%
\theoriaref{Θεώρημα}{#1}{Η ζητούμενη απόδειξη δίνεται στην αντίστοιχη απόδειξη.}}

\begin{thema}{Α}
\begin{erwthma}
\item \apodeixiref{thewrhma:paragwgos_ekthetikhs}

\item \orismosref{orismos:grafikh_parastash}

\item \orismosref{orismos:rythmos_metavolhs}

\item Οι χαρακτηρισμοί των προτάσεων είναι:
\begin{alist}
\item \textbf{Σωστό}. Αν οι $f,g$ είναι παραγωγίσιμες, τότε
\[ (f(x)+g(x))'=f'(x)+g'(x). \]
\item \textbf{Λάθος}. Το εμβαδόν δίνεται γενικά από
\[ E=\dint_a^\beta |f(x)|\d x, \]
όχι πάντα από το $\dint_a^\beta f(x)\d x$.
\item \textbf{Σωστό}. Ο κανόνας \eng{De L'Hospital} εφαρμόζεται μόνο σε απροσδιοριστίες της μορφής $\frac{0}{0}$ ή $\frac{\pm\infty}{\pm\infty}$, όταν ισχύουν οι προϋποθέσεις του.
\item \textbf{Σωστό}. Για κυρτή συνάρτηση η γραφική παράσταση βρίσκεται πάνω από κάθε εφαπτομένη της, άρα
\[ f(x)\ge f(x_0)+f'(x_0)(x-x_0). \]
\item \textbf{Σωστό}. Αν ένα πολυώνυμο βαθμού μεγαλύτερου του $2$ έχει δύο τοπικά ακρότατα στα $x_1<x_2$, τότε από το θεώρημα \eng{Rolle} η $f'$ έχει δύο ρίζες και η $f''$ μηδενίζεται σε κάποιο σημείο του $(x_1,x_2)$, όπου αλλάζει η μονοτονία της $f'$. Έτσι υπάρχει σημείο καμπής.
\end{alist}

\item Σωστή απάντηση είναι η \textbf{β}. Αν $\lim\limits_{x\to x_0}f(x)=-\infty$, τότε για τιμές του $x$ αρκετά κοντά στο $x_0$ ισχύει $f(x)<0$.
\end{erwthma}
\end{thema}
%=========== ΤΕΛΟΣ ΛΥΣΗΣ - 22ο ΘΕΜΑ Α ===========
